\documentclass{article}
\usepackage[T1]{fontenc}
\usepackage{lmodern}
\usepackage{graphicx}
\usepackage{amsmath,amssymb}
\usepackage[a4paper,margin=2cm]{geometry}
\usepackage[absolute,overlay]{textpos}
\setlength{\TPHorizModule}{1mm}
\setlength{\TPVertModule}{1mm}
\usepackage[indonesian]{babel}
\usepackage{hyperref}
\setlength{\emergencystretch}{2em}
% unit: Halaman 32

\begin{document}

% === Halaman 32 (llm) ===

\section*{Penggandaan Titik di dalam EC pada GF(p)}

Misalkan $P(x_p, y_p)$ yang dalam hal ini $y_p \neq 0$.
Penggandaan titik: $2P = R$

Koordinat Titik R:
\begin{align*}
x_r &\equiv m^2 - 2x_p \pmod{p} \\
y_r &\equiv m(x_p - x_r) - y_p \pmod{p}
\end{align*}

Yang dalam hal ini,
\[ m \equiv \frac{3x_p^2 + a}{2y_p} \pmod{p} \]

Jika $y_p = 0$ maka $m$ tidak terdefinisi sehingga $2P = O$

\end{document}
